\documentclass[12pt]{article}
\usepackage[T1]{fontenc}
\usepackage[utf8]{inputenc}
\usepackage{lmodern}
\usepackage{textcomp}
\usepackage{enumitem}
\usepackage{geometry}
\geometry{margin=1in}
\begin{document}
\begin{center}
Answer Key - Health Sciences - Graduate Level Assessment 8 \
\small (Answers are ordered exactly as in the question paper.)
\end{center}

\section*{Multiple Choice}
\begin{enumerate}[label=\arabic*.]
\item \textbf{Answer:} A
\\ \textit{Options:} A) World Medical Association (WMA); B) United Nations Educational, Scientific and Cultural Organization (UNESCO); C) Geneva Convention; D) World Health Organization (WHO); E) Declaration of Helsinki
\\ \textit{Note:} The World Medical Association (WMA) established the Declaration of Helsinki, which outlines the fundamental principles of ethical research, including respect for persons, beneficence, and justice.
\item \textbf{Answer:} A
\\ \textit{Options:} A) More than 90\%; B) Less than 10\%; C) More than 50\%; D) Less than 80\%; E) Exactly 50\%
\\ \textit{Note:} The correct answer is supported by research indicating that over 90\% of adults in various populations engage in some form of regular physical activity, underscoring the widespread recognition of the importance of exercise for health.
\item \textbf{Answer:} A
\\ \textit{Options:} A) Alcohol.; B) Synthetic cannabinoids.; C) Cocaine.; D) Benzodiazepines.; E) Prescription painkillers.
\\ \textit{Note:} Alcohol is the most commonly responsible substance for acute confusion in patients due to its widespread use and the potential for intoxication or withdrawal effects that significantly impair cognitive function.
\item \textbf{Answer:} A
\\ \textit{Options:} A) Eco-friendly; B) Sunshine; C) Forest; D) Rapid decay; E) Organic
\\ \textit{Note:} An eco-friendly cemetery facilitates natural decomposition by allowing bodies to return to the earth without the use of caskets, aligning with sustainable practices that minimize environmental impact.
\item \textbf{Answer:} C
\\ \textit{Options:} A) Loss of muscle mass; B) Shrinkage of the thymus gland; C) Differences in weight; D) Changes in metabolism rate; E) Changes in eyesight
\\ \textit{Note:} Differences in weight could reflect variations between different generational cohorts due to factors such as diet, lifestyle, and environmental influences, rather than representing a universal change that occurs with aging.
\end{enumerate}

\section*{True / False}
\begin{enumerate}[label=\arabic*.]
\setcounter{enumi}{5}
\item \textbf{Answer:} False
\\ \textit{Options:} A) True; B) False
\\ \textit{Note:} Antibiotics are not universally recommended prior to catheterization for patients with skin infections, as the decision is influenced by the patient's overall medical history, including potential risks and the specific type of infection.
\item \textbf{Answer:} True
\\ \textit{Options:} A) True; B) False
\\ \textit{Note:} The morphology of the Dane particle, which is the complete virus particle of Hepatitis B, is indeed characterized by its unique flexuous and filamentous structure, distinguishing it from other viral forms.
\item \textbf{Answer:} False
\\ \textit{Options:} A) True; B) False
\\ \textit{Note:} The statement is false because current estimates suggest that the number of human protein-coding genes is approximately 20,000 to 25,000, significantly lower than the claimed range of 50,000 to 60,000.
\item \textbf{Answer:} True
\\ \textit{Options:} A) True; B) False
\\ \textit{Note:} In hypovolaemic shock, the body can compensate for a loss of up to 15\% of blood volume by increasing peripheral resistance and maintaining cardiac output, which allows heart rate and blood pressure to remain relatively stable until the compensatory mechanisms are overwhelmed.
\item \textbf{Answer:} False
\\ \textit{Options:} A) True; B) False
\\ \textit{Note:} Urbanization and rising incomes often lead to changes in dietary patterns that include increased consumption of grains and pulses due to their nutritional value and affordability, rather than a decline in their consumption.
\end{enumerate}

\section*{Long Form Response}
\begin{enumerate}[label=\arabic*.]
\setcounter{enumi}{10}
\item \textbf{Answer:} Fat digestion and absorption are critical processes that hinge on the emulsification and digestion of triglycerides, or triacylglycerols. Initially, dietary fats undergo emulsification in the small intestine, where bile salts facilitate the breakdown of large fat globules into smaller droplets, increasing the surface area for enzymatic action. This is followed by the hydrolysis of triglycerides by pancreatic lipase, which cleaves them into free fatty acids and monoglycerides. Once digested, these products aggregate with bile salts to form micelles, which transport the fatty acids across the intestinal epithelium for absorption. The efficient emulsification and subsequent digestion of triglycerides are essential for the transportation of fatty acids, ensuring their availability for cellular metabolism and energy production.
\\ \textit{Note:} The answer correctly outlines the essential steps of fat digestion and absorption by detailing how emulsification by bile salts enhances the enzymatic hydrolysis of triglycerides by pancreatic lipase, leading to the production of free fatty acids and monoglycerides, which are crucial for the transportation of fatty acids into the body.
\item \textbf{Answer:} White bread is primarily composed of refined wheat flour, which undergoes processing that removes the bran and germ, leaving mainly the starchy endosperm. This processing results in a product that is low in fiber and certain micronutrients, even though it can provide a quick source of carbohydrates. Notably, white bread is classified as not containing free sugars because it does not have added sugars or sugar substitutes; its carbohydrate content mainly consists of complex carbohydrates that break down into glucose during digestion. The nutritional implications of consuming white bread include a rapid spike in blood sugar levels due to its high glycemic index, potentially leading to negative health effects such as increased risk of type 2 diabetes and cardiovascular disease. Therefore, while white bread serves as a staple food, its lack of fiber and high glycemic response highlight the importance of considering whole grain alternatives for better health outcomes.
\\ \textit{Note:} correct because it accurately describes the composition of white bread as primarily consisting of refined wheat flour with complex carbohydrates, clarifying that it lacks added sugars or sugar substitutes, which is why it is classified as not containing free sugars.
\end{enumerate}

\vfill
\noindent\textit{End of Answer Key}
\end{document}